% Checked auxiliary and conditional CM results. See cm_transport.md and CHECK_EVENTS.tsv.
\section{Auxiliary CM transport and a conditional secant construction}
\label{sec:cm-transport}

The linear and descent results in this section are independent of a
conjectural cycle construction. A separately labelled theorem uses
Markman's unrefereed source sheaf conditionally. The explicit
rank-zero construction elsewhere in this paper supersedes that
conditional route for the balanced Gaussian CM example; the conditional
route is retained with its precise scope and its unsuccessful boundary
specialization recorded.

\begin{proposition}[SF-CM01: semiregularity under a fixed isogeny]
\label{SF-CM01}
Let $f:A\to Y$ be an isogeny of complex abelian varieties and let
$\overline F$ be an ordinary semiregular coherent sheaf on $Y$.
Put $G=\ker f$ and $F=f^*\overline F$. Then $F$ is a perfect
coherent sheaf with no negative self-Ext and is $G$-semiregular for
the canonical coherent action of the actual finite subgroup
$G\subset A\times\{0\}$ by translations. If its rank is $r\ne0$,
then
\[
 e^{-c_1(F)/r}\operatorname{ch}(F)
 =f^*\bigl(e^{-c_1(\overline F)/r}\operatorname{ch}(\overline F)\bigr).
\]
\end{proposition}
\begin{proof}
The isogeny is finite etale. Its translations satisfy $f\circ t_g=f$,
so functorial pullback gives canonical isomorphisms $t_g^*F\simeq F$
satisfying the cocycle identity. This is an honest coherent action,
with no projective ambiguity or invisible central kernel. Finite
etale descent, and exactness of taking invariants in characteristic
zero, give
\[
 \operatorname{Ext}^2_A(F,F)^G
 \simeq\operatorname{Ext}^2_Y(\overline F,\overline F).
\]
The Atiyah class and its trace commute with etale pullback. Thus the
full semiregularity map restricted to this invariant subspace is
$f^*\circ\sigma_{\overline F}$. On every Hodge-cohomology summand
of $\bigoplus_qH^{q+2}(Y,\Omega_Y^q)$ the pullback is injective,
since $f_*f^*=(\deg f)\operatorname{id}$. The restricted map is
therefore injective. This is strong $G$-semiregularity, and hence
weak $G$-semiregularity in the convention of DEG\_L S-14; see also
\cite[Definition~2.6 and Lemma~2.7, pp.~5--6]{SF-Perry2604}.
Every coherent sheaf on the smooth $A$ is perfect, and negative
self-Ext vanishes for sheaves in degree zero. Naturality of rank and
Chern character proves the final equality.
\end{proof}
This statement uses the full Atiyah-trace map. Injectivity of a
single middle component is not being substituted for semiregularity.

\begin{proposition}[SF-CM02: the entire invariant character]
\label{SF-CM02}
Let $(A,K,h)$ be a polarized Weil sixfold of signature $(3,3)$.
If a rational even cohomology class $\kappa$ remains Hodge under
all local deformations in its full nine-dimensional Weil period
domain, then
\[
 \kappa=P(h)+t,\qquad P\in\mathbf Q[x],\quad t\in W_K\subset H^6(A,\mathbf Q).
\]
If $\kappa_3$ is not proportional to $h^3$, then $t\ne0$.
\end{proposition}
\begin{proof}
The Hodge circle of a balanced Weil point lies in $\operatorname{SU}(3,3)$:
on either six-dimensional $K$-eigenspace its determinant is
$z^3z^{-3}=1$. Its conjugates are the Hodge circles in the period
domain. Analyticity extends the assumed Hodge condition from an open
neighborhood to the connected period domain. The infinitesimal
generators of these circles generate a nonzero ideal in the simple
real Lie algebra $\mathfrak{su}(3,3)$, hence all of that Lie algebra.
The class $\kappa$ is consequently $\operatorname{SU}(3,3)$-invariant.

After scalar extension to $\mathbf C$, write $H^1=U\oplus U^*$,
where $U$ is the standard six-dimensional representation of
$\operatorname{SL}_6$. The decomposition
\[
 \bigwedge^k(U\oplus U^*)
 =\bigoplus_{a+b=k}\bigwedge^aU\otimes\bigwedge^bU^*
\]
contains an invariant exactly when $a=b$, or when
$(a,b)=(6,0),(0,6)$. Indeed the exterior powers of the standard
representation are irreducible with distinct highest weights, except
that $\bigwedge^0U$ and $\bigwedge^6U$ are both trivial. For $a=b$
the invariant line is generated by the $a$th power of the invariant
pairing $h$. The two exceptional lines are the determinant lines
$\bigwedge^6U$ and $\bigwedge^6U^*$ in degree six. Thus the entire
invariant subspace is
\[
 \mathbf Q[h]/(h^7)\oplus W_K.
\]
The polynomial summand already includes $\mathbf Qh^3$. The
decomposition is rational, since $h$ and the sum of the conjugate
determinant lines are rational. The asserted formula and nonvanishing
criterion follow.
\end{proof}
This supplies the full-character conclusion needed by S-14, rather
than a conclusion about the middle character alone.

\begin{proposition}[SF-CM04: the fixed hyperbolic spin model forces splitness]
\label{SF-CM04}
Let $(V,q)$ be a hyperbolic rational quadratic space of dimension
$4n$. Suppose $f^2=-d$, $d>0$, and
$q(fx,y)=-q(x,fy)$. Endow $V$ with its
$K=\mathbf Q(\sqrt{-d})$ structure through $f$. The Hermitian form
\[
 H(x,y)=d\,q(x,y)+\sqrt{-d}\,q(fx,y)
\]
(up to interchanging its linear argument convention) is split of
$K$-rank $2n$. Its determinant class is $[(-1)^n]$.
\end{proposition}
\begin{proof}
Skew-adjointness gives $q(fx,x)=0$, hence $H(x,x)=dq(x,x)$.
A nonzero rational isotropic vector for the hyperbolic $q$ is
therefore $K$-isotropic for $H$. Nondegeneracy splits off a
Hermitian hyperbolic plane containing it. The rational quadratic
form $\operatorname{Re}H/d$ on this plane is hyperbolic of rank
four: the two isotropic $K$-lines are two perfectly paired isotropic
rational planes. The Hermitian orthogonal complement is also
$q$-orthogonal and $K$-stable. Rational Witt cancellation makes
$q$ on that complement hyperbolic. Repeating gives $n$ Hermitian
hyperbolic planes, with determinant $(-1)^n$ modulo norms.
\end{proof}
In Markman's construction the ambient lattice is
$H^1(X)\oplus H^1(\widehat X)$ with the evaluation pairing, a sum
of $2n$ hyperbolic planes. His oriented $K$-secant endomorphism
and form satisfy these identities
\cite[\S3.1, equation~(3.1.2) and Lemma~3.1.2, pp.~22--23]{SF-Markman2502}.
Consequently the split determinant is not restricted to the explicit
ample secants $\exp(\pm\sqrt{-d}\Theta)$. Another secant in this
same quadratic model cannot repair the nonsplit determinant gap.
The linear-algebra proposition is unconditional; geometric uses of
the unrefereed construction remain conditional.

\begin{proposition}[SF-CM05: the literal elliptic-chain boundary fails]
\label{SF-CM05}
On $A_3=E_1\times E_2\times E_3$, for $b\in E_2\setminus\{0\}$,
put
\[
 C_b=(E_1\times\{0\}\times\{0\})
 \cup(\{0\}\times E_2\times\{0\})
 \cup(\{0\}\times\{b\}\times E_3).
\]
Its ideal sheaf is not semiregular. If a finite translation subgroup
$H\subset A_3$ gives pairwise disjoint translates of $C_b$, the
ideal of their union with its canonical $H$-linearization is not
weakly $H$-semiregular. The defect persists for the natural product
action on exterior products of two such orbit ideals, and under line
bundle twists or derived equivalences.
\end{proposition}
\begin{proof}
The curve is a connected genus-three chain with two ordinary nodes.
At either node its ideal has local form $(z,xy)$ in the smooth
threefold, so it is a local complete intersection. Its normal sheaf
has at least four independent global sections: the three translations
of $A_3$, and the variation of $b$. To verify independence, restrict
to the smooth portions of the three components. A translation vector
$(v_1,v_2,v_3)$ has normal components
\[
 (v_2,v_3),\qquad(v_1,v_3),\qquad(v_1,v_2).
\]
The $b$-variation is zero on the first two components and nonzero
only in the $E_2$-normal direction of the third. A translation
matching its first two restrictions must be zero and cannot match
the third. The same calculation shows that the three translation
directions themselves are independent. These are Hilbert tangent
vectors: $C_b$ is a flat family locally in $b$, with constant nodal
local equations, component genera and incidence relations.

Write $I=\mathcal I_{C_b}$. In the exact sequence obtained by applying
$\operatorname{Hom}(I,-)$ to
$0\to I\to\mathcal O_{A_3}\to\mathcal O_{C_b}\to0$,
the map $\operatorname{Hom}(I,I)\to\operatorname{Hom}(I,\mathcal O_{A_3})$
is the isomorphism $\mathbf C\to\mathbf C$. Consequently
\[
 H^0(N_{C_b/A_3})=\operatorname{Hom}(I,\mathcal O_{C_b})
 \hookrightarrow\operatorname{Ext}^1(I,I).
\]
Tensoring by $\operatorname{Pic}^0(A_3)$ gives another
three-dimensional subspace $H^1(\mathcal O_{A_3})$. Its trace is
the identity because $I$ has rank one, whereas the embedded ideal
deformations have trivial determinant. Thus the two subspaces are
independent and $\operatorname{ext}^1(I,I)\ge7$.
Serre duality on the abelian threefold gives
$\operatorname{ext}^2(I,I)\ge7$. The full semiregularity target has
dimension
\[
 h^2(\mathcal O_{A_3})+h^3(\Omega^1_{A_3})=3+3=6,
\]
so the map is not injective.

For the orbit union $Z_b=\bigcup_{h\in H}(C_b+h)$, disjointness
persists locally in $b$. Its invariant normal directions contain the
same four-dimensional diagonal subspace. The three trace directions
from $H^1(\mathcal O_{A_3})$ are invariant. Translation preserves
the volume form, so equivariant Serre duality yields
$\dim\operatorname{Ext}^2(\mathcal I_{Z_b},\mathcal I_{Z_b})^H\ge7$.
More decisively, descent identifies the invariant object with the
ideal of the embedded image of $C_b$ in the abelian threefold
$A_3/H$. That object's ordinary semiregularity target again has
dimension six. It is therefore not semiregular in the invariant
category, which is exactly failure of weak $H$-semiregularity.
The automorphisms of the ideal are scalars; other linearizations
only tensor its descent by a torsion line bundle and do not repair
the kernel.

On the exterior product, take a nonzero kernel element in one
factor's $\operatorname{Ext}^2$ and tensor it with the other
factor's identity. The Kunneth summand is nonzero but its full trace
is zero. Twisting by line bundles multiplies the trace by an
invertible exponential; derived equivalence transports the
categorical trace and its kernel. Thus the defect persists under
the stated operations.
\end{proof}
This proposition concerns the literal boundary chain and its natural
orbit repair. It does not obstruct smooth CM Jacobians or arbitrary
complexes on $E^6$. In particular, the chain used in an openness
argument cannot itself be substituted for the required semiregular
CM seed without an additional proof.

\subsection{Conditional subgroup: Markman's source sheaf}
\label{sec:conditional-cm-source}

\begin{theorem}[SF-CM03: a conditional balanced CM seed]
\label{SF-CM03}
Assume the sheaf construction and semiregularity assertions of
Markman \cite{SF-Markman2502}, specifically Theorem~1.4.1(2)--(4),
Lemmas~9.1.4, 9.3.1, 9.3.5, 9.3.11 and Remark~9.3.7.
Let $K=\mathbf Q(\sqrt{-m})$, $m>0$ squarefree, and let $E$ be
any elliptic curve with $\operatorname{End}_{\mathbf Q}(E)=K$.
On $A=E^6$ with the standard balanced action
$\operatorname{diag}(a,a,a,\bar a,\bar a,\bar a)$, there exist:
an integral ample compatible polarization $\theta$ of determinant
$[-1]$, a reflexive coherent sheaf $F$ of rank $32m$, and an honest
finite translation group $G\subset A\times\{0\}$ for which $F$
is $G$-semiregular and has no negative self-Ext, such that
\[
 e^{B_0}\operatorname{ch}(F)=P(\theta)+t,
 \quad B_0=-c_1(F)/(32m),\quad P\in\mathbf Q[x],\quad0\ne t\in W_K.
\]
The equality holds in the full even cohomology. This statement has no
class-number hypothesis. It does not assert a prescribed principal
polarization, a prescribed Weil direction, or a numerical detector
bound.
\end{theorem}
\begin{proof}
We specify the finite choices and identify precisely where the
conditional source assertions are used.

\emph{The radicand.} Set $d=4m$, $N=d+1$, and $r=8d=32m$.
Then $d\ge4$ is even, $K=\mathbf Q(\sqrt{-d})$, and
$\gcd(r,N^2)=1$. Thus the construction satisfies the source
hypothesis $d\ge3$ and admits its untwisted determinant-power
normalization. In particular the inadmissible $d=1$ Gaussian
specialization is avoided.

\emph{A good smooth CM Jacobian.} Let $\tau_0\in\mathfrak H_3$
be the period matrix of the product principally polarized $E^3$.
The orbit $\operatorname{Sp}_6(\mathbf Q)\tau_0$ is analytically
dense. Indeed the Cayley map
$X\mapsto(1-X)(1+X)^{-1}$ on the symplectic Lie algebra is rational
and parametrizes a neighborhood of the identity. Rational
Lie-algebra points are dense there. As
$\operatorname{Sp}_6(\mathbf R)$ is connected, its rational points
are dense throughout, and it acts transitively on $\mathfrak H_3$.
A rational symplectic transform changes the marked lattice to a
commensurable one, so every orbit point is a principally polarized
abelian variety isogenous to $E^3$.

The locus of nonhyperelliptic genus-three Jacobians contains an
analytic open subset of $\mathcal A_3$. To check the needed rank,
a smooth plane quartic has $K_C=\mathcal O_C(1)$, and
\[
 \operatorname{Sym}^2H^0(K_C)\longrightarrow H^0(K_C^2)
\]
is an isomorphism between six-dimensional spaces: no nonzero
quadric contains the quartic. This is the dual period differential,
so the differential has rank six. The good orbit and intersection
hypotheses in the conditional source Lemmas~9.1.4 and~9.3.1 hold
on a nonempty open subset of the finite moduli cover recording the
$N$-torsion groups. A finite etale cover is open, so the image gives
a nonempty analytic open subset of good Jacobians in $\mathcal A_3$.
Density supplies $q\in\operatorname{Sp}_6(\mathbf Q)$ in that
open subset. Choose its smooth curve $C$, put $X=J(C)$, and choose
the good cyclic groups $G_1,G_2$ of order $N$. Then $X$ is isogenous
to $E^3$.

This choice establishes existence in a nonempty open subset. No
specific rational matrix, plane quartic equation or torsion point
coordinates are asserted to have been computed for this conditional
construction.

\emph{An ordinary semiregular quotient.} The source theorem gives
a simple reflexive sheaf $\mathcal E$ of rank $r$ on
$M=X\times\widehat X$, together with finite translation/tensor
symmetries. Lemma~9.3.5, p.~84, uses even $d$ and
$\gcd(r,N^2)=1$ to cancel the tensor part by a determinant power.
One may take an integer satisfying $ar\equiv-1\pmod{N^2}$,
with the sign adjusted to the source convention for the intermediate
derived object. The result descends along a translation quotient
$q_M:M\to Y$ to an ordinary reflexive sheaf $\overline F$ of
rank $r$. Remark~9.3.7, p.~86, identifies the normalized twisted
sheaf $\mathcal B$ with the normalization of this ordinary sheaf;
Lemma~9.3.11, p.~88, states that $\mathcal B$ is semiregular.

This also proves ordinary semiregularity of $\overline F$.
Replacing its Atiyah class by
$\operatorname{at}_{\overline F}-(c_1(\overline F)/r)\operatorname{id}$
multiplies the entire Atiyah-trace map by
$\exp(-c_1(\overline F)/r)$. On the full target
$\bigoplus_qH^{q+2}(\Omega_Y^q)$ this is an invertible triangular
operator. The normalized local scalar twists cancel in the
endomorphism complex, so the two self-Ext spaces agree. Thus
injectivity of the normalized twisted map is equivalent to
injectivity of the ordinary map. Moreover
\[
 q_M^*\kappa(\overline F)=\kappa(\mathcal E),\qquad
 \kappa(F')=e^{-c_1(F')/\operatorname{rank}(F')}\operatorname{ch}(F').
\]

\emph{The balanced action on the actual power.} Since $Y$ is
isogenous to $X\times\widehat X$, it is isogenous to $E^6$ and
$\operatorname{End}_{\mathbf Q}(Y)\simeq M_6(K)$. Here the central
copy of $K$ initially comes from the elliptic CM action. Markman's
Weil embedding $\eta:K\to M_6(K)$ commutes with this center.
Over the central $K$ its generator has two eigenvalues
$\pm\sqrt{-d}$. The two eigenspaces have $K$-dimension three,
because their holomorphic dimensions are the Weil multiplicities
$(3,3)$. A matrix in $\operatorname{GL}_6(K)$ therefore conjugates
$\eta$ to the specified balanced diagonal embedding. Compose with
an isogeny identification and clear denominators to obtain a
$K$-linear isogeny $f:E^6_{\rm balanced}\to Y$. Equality of the
elliptic CM field and the secant field is essential to this
argument; it is not claimed for an unrelated elliptic field.

\emph{Pullback and the entire character.} Set $F=f^*\overline F$
and $G=\ker f$. SF-CM01 gives the honest translation action,
strong $G$-semiregularity and negative-Ext vanishing. The positive
rational Weil polarization $h_Y$ pulls back to a compatible
positive class. Choose $c>0$ such that $\theta=cf^*h_Y$ is an
integral ample line-bundle class. The source determinant is $[-1]$
by its Lemma~3.1.3, pp.~23--24, or by SF-CM04. Quotient and
$K$-isogeny transport preserve it, and scaling by $c$ changes the
rank-six determinant by $c^6=\operatorname{Nm}(c^3)$, leaving its
class unchanged.

Theorem~1.4.1(3) states that the entire $\kappa(\mathcal E)$ stays
Hodge on the full Weil period domain. Fixed isogenies identify the
local period domains, so the same is true of $\kappa(F)$.
SF-CM02 gives $\kappa(F)=P(\theta)+t$. By
Theorem~1.4.1(4), the source $\kappa_3(\mathcal E)$ has nonzero
Weil projection: its $K$-translates together with $h^3$ span three
dimensions. Isogeny pullback is injective and preserves the Weil
plane, so $t\ne0$. Finally $B_0=-c_1(F)/r$ is algebraic and
$e^{B_0}\operatorname{ch}(F)=\kappa(F)$, as required. Replacing
$\theta$ by a very ample positive multiple merely rescales $P$.
\end{proof}

\paragraph{What the conditional theorem changes.}
It constructs a genuine coherent object at the balanced CM point,
subject to the stated preprint input. The first-order assertion of
Theorem~1.4.1 alone would not suffice: the finite quotient and
Lemma~9.3.11 are essential here. Choosing even $d$ removes the gerbe
coherence issue entirely. SF-CM02 supplies the whole character, not
just its middle term. One nonzero Weil direction is enough for the
detector programme; its rational $K$-translates span the plane.
There is no assertion that arbitrary direct sums preserve
semiregularity while prescribing another direction.

Applying the supplied DEG\_L S-14/Perry theorem gives bounded
division on the associated split determinant family conditionally
on these source assertions. The resulting bound remains
existential until the chosen CM Jacobian, sheaf presentation and
degree data are made explicit. It is not presented as a numerical
detector constant, and the chosen polarization on $E^6$ is not
identified with its product principal polarization. Nonsplit
determinants remain outside this fixed hyperbolic secant model.

In the programme's language, this is a conditional $K$-theory input
to S-14, with a full $\lambda$-ring Chern character and an object
whose deformation theory is controlled. The subsequent uniform
integral-cycle step uses the supplied HECKE\_DIV/GHOSTS reduction,
condition A at cote~16 scan~050. Neither condition I (scan~049 and
cote~44 Theorem~4) nor a bound on the CM isogeny or core index $J$
produces the required object. No additional archive or EGA/SGA
locator is claimed for the new representation and descent arguments.
